% Generated draft: requires mathematical and source review.
\documentclass[11pt]{article}
\usepackage[margin=25mm]{geometry}
\usepackage{fontspec}
\setmainfont{Latin Modern Roman}
\usepackage{amsmath,amssymb,mathtools,unicode-math}
\setmathfont{Latin Modern Math}
\usepackage{enumitem,needspace,xurl,hyperref}
\setlength{\parindent}{0pt}
\setlength{\parskip}{5pt}
\setlength{\emergencystretch}{3em}
\newcommand{\N}{\mathbb{N}}
\newcommand{\Z}{\mathbb{Z}}
\newcommand{\Q}{\mathbb{Q}}
\newcommand{\R}{\mathbb{R}}
\newcommand{\C}{\mathbb{C}}
\newcommand{\F}{\mathbb{F}}
\newcommand{\A}{\mathbb{A}}
\newcommand{\PP}{\mathbb{P}}
\newcommand{\E}{\mathbb{E}}
\newcommand{\eps}{\varepsilon}
\newcommand{\dd}{\,\mathrm{d}}
\begin{document}
\section*{UnsolvedMath ID: 800370}
\textbf{Source problem code:} catalogue entry
\subsection{Definitions and mathematical statement}
For a positive integer \(n\), let \([n]=\{1,\ldots,n\}\), and let \(S_n\) be the set of all permutations of \([n]\). For \(\sigma\in S_n\), define its \emph{odd length}, also called its \emph{odd inversion number}, by
\[
L(\sigma)=\left|\left\{(i,j)\in[n]^2:i<j,\ \sigma(i)>\sigma(j),\ i\not\equiv j\pmod 2\right\}\right|.
\]
Define
\[
L_n(x)=\sum_{\sigma\in S_n}x^{L(\sigma)}
      =\sum_{k=0}^{d_n}a_{n,k}x^k,
\qquad d_n=\deg L_n(x),
\]
where \(a_{n,k}=|\{\sigma\in S_n:L(\sigma)=k\}|\).

The supplied conjecture states that, for every \(n\in\mathbb{P}\) with \(n\geq5\), the polynomial \(L_n(x)\) is symmetric and unimodal. Explicitly, the required properties are
\[
a_{n,k}=a_{n,d_n-k}\qquad(0\leq k\leq d_n)
\]
and the existence of an integer \(m\in\{0,\ldots,d_n\}\) such that
\[
a_{n,0}\leq\cdots\leq a_{n,m}\geq\cdots\geq a_{n,d_n}.
\]
The supplied record does not define \(\mathbb{P}\). Its notation is retained here: whether the quantifier ranges over all positive integers or only primes requires clarification from the source. The record labels this assertion open; that status is supplied metadata, not an independently verified mathematical assessment.
\subsection{Short English statement}
Count inversions whose two positions have opposite parity. The conjecture asserts that, for each permitted \(n\geq5\), the numbers of permutations having each possible odd inversion number form a symmetric, unimodal sequence. The source's domain notation \(\mathbb{P}\) remains unresolved.
\subsection{Sources}
\begin{itemize}
\item Local UnsolvedMath registry, permanent ID \(800370\), exact source problem code \texttt{catalogue entry}, titled ``Coxeter Groups and Unimodality --- Conjecture 8.'' The record attributes the problem to Brenti (2024) and supplies the status \texttt{open}.
\item Brenti, \emph{Some Open Problems on Coxeter Groups and Unimodality} (2024), Conjecture 8, \url{https://arxiv.org/abs/2410.09897}. The supplied provenance records access on 2026-10-16 and extraction by \texttt{arxiv\_latex}; these are supplied provenance details, not a claim of independent source access.
\item The supplied statement identifies the conjecture as [BreCa2, Conj. 6.2]. It also reports that the statistic was introduced in [KloVol] in relation to formed spaces and was further studied in [BreCav] and [BreCa2]. Complete bibliographic details for these reference labels are not supplied.
\end{itemize}
\end{document}
